\documentclass[11pt]{article}
\usepackage{shared/template_tgir_article}
\usepackage{booktabs}
\usepackage{longtable}
\usepackage{tabularx}
\usepackage{array}
\usepackage{enumitem}
\usepackage{xcolor}
\usepackage{listings}

\lstset{
    basicstyle=\ttfamily\small,
    breaklines=true,
    columns=fullflexible,
    frame=single
}

\newcommand{\product}{TGIR QuantLib Tools}
\newcommand{\code}[1]{\texttt{#1}}
\newcommand{\docdate}{2026-04-06}

% Visual design is controlled exclusively by template_tgir_article.sty and
% IEEEtran. Keep this file limited to content-support packages and commands.


\title{Deployment, CI/CD, And DigitalOcean Guide\\\large \product}
\author{TG Investments and Research}
\date{\docdate}

\begin{document}
\maketitle
\tableofcontents
\newpage

\section{Purpose}
This guide describes a practical production path for the repo using GitHub Actions for CI/CD and a single DigitalOcean droplet for runtime hosting.

\section{Target Architecture}
The recommended production layout is:
\begin{enumerate}[leftmargin=1.5em]
    \item GitHub repository for source control and CI/CD,
    \item one Ubuntu-based DigitalOcean droplet,
    \item Apache as the public HTTPS reverse proxy,
    \item one local WSGI process managed by systemd,
    \item a release directory layout with a \code{current} symlink, and
    \item a shared server-side \code{.env} file that is never committed.
\end{enumerate}

\section{Files Added For Deployment}
This repo now contains deployment-supporting assets:
\begin{itemize}[leftmargin=1.5em]
    \item \path{wsgi.py},
    \item \path{requirements-production.txt},
    \item \path{deploy/systemd/tgir-quantlib-tools.service.template},
    \item \path{deploy/apache/tgir-quantlib-tools.conf.template},
    \item \path{deploy/env/production.env.example},
    \item \path{.github/workflows/ci.yml}, and
    \item \path{.github/workflows/deploy_digitalocean.yml}.
\end{itemize}

\section{CI Pipeline}
\subsection{Objective}
CI should prove that the code installs, compiles, the LaTeX docs build, and the quantitative tests pass before deployment is allowed.

\subsection{Recommended CI Stages}
\begin{enumerate}[leftmargin=1.5em]
    \item checkout and Python setup,
    \item virtual-environment creation,
    \item dependency installation,
    \item Python syntax compilation,
    \item targeted unittest execution,
    \item LaTeX PDF build for the doc set, and
    \item artifact upload of generated PDFs.
\end{enumerate}

\section{CD Pipeline}
\subsection{One-Time Bootstrap}
Before automated deployment, bootstrap the droplet manually:
\begin{enumerate}[leftmargin=1.5em]
    \item create the application root such as \path{/opt/tgir_quantlib_tools},
    \item create a shared \code{.env} file under \path{/opt/tgir_quantlib_tools/shared/.env},
    \item install Apache and enable proxy-related modules,
    \item install the systemd service using the template in \path{deploy/systemd/},
    \item install the Apache site configuration using the template in \path{deploy/apache/}, and
    \item confirm that \code{/health} works on the local WSGI bind port and through Apache.
\end{enumerate}

\subsection{Automated Release Flow}
The deployment workflow should:
\begin{enumerate}[leftmargin=1.5em]
    \item run after CI succeeds on \code{main} or manual dispatch,
    \item upload the repository as a release tarball,
    \item create a timestamped or commit-hash release directory on the droplet,
    \item create a fresh \code{.venv} inside that release,
    \item install \path{requirements-production.txt},
    \item link the shared \code{.env} and debug directories into the new release,
    \item repoint \code{current} to the new release,
    \item restart the systemd service, and
    \item smoke-check \code{/health}.
\end{enumerate}

\section{Droplet Preparation}
\subsection{Packages}
On Ubuntu, the typical bootstrap is:
\begin{lstlisting}
sudo apt-get update
sudo apt-get install -y python3 python3-venv python3-pip apache2
\end{lstlisting}

\subsection{Apache Modules}
Enable the required modules:
\begin{lstlisting}
sudo a2enmod proxy proxy_http headers rewrite ssl
\end{lstlisting}

\subsection{Filesystem Layout}
Recommended layout:
\begin{lstlisting}
/opt/tgir_quantlib_tools/
  current -> /opt/tgir_quantlib_tools/releases/<release-id>
  releases/
  shared/
    .env
    debug/
\end{lstlisting}

\section{Production Configuration}
\subsection{WSGI Entry Point}
Use \path{wsgi.py} as the production entrypoint so the systemd service can run a local WSGI server cleanly.

\subsection{Production Dependencies}
Use \path{requirements-production.txt}. It extends the development requirements with \code{gunicorn} for the production process manager.

\subsection{Service Template}
The systemd template starts \code{gunicorn} on a loopback port such as \code{127.0.0.1:8008}. Apache proxies public HTTPS requests to that local port.

\section{GitHub Secrets}
The deployment workflow assumes these secrets or environment variables:
\begin{tabularx}{\textwidth}{>{\raggedright\arraybackslash}p{0.32\textwidth}X}
\toprule
Secret & Purpose \\
\midrule
\code{DO\_HOST} & Droplet hostname or IP. \\
\code{DO\_USER} & SSH user for deployments. \\
\code{DO\_SSH\_KEY} & Private key used by GitHub Actions. \\
\code{DO\_APP\_ROOT} & App root, for example \path{/opt/tgir_quantlib_tools}. \\
\code{DO\_HEALTHCHECK\_URL} & Public health endpoint for final verification. \\
\bottomrule
\end{tabularx}

\section{Suggested Release Procedure}
\begin{enumerate}[leftmargin=1.5em]
    \item Open a pull request and let CI pass.
    \item Merge to \code{main}.
    \item Let the deployment workflow publish the new release to the droplet.
    \item Verify \code{/health}, login, dashboard rendering, and the cliquet page.
    \item If needed, roll back by repointing the \code{current} symlink to the previous release and restarting the service.
\end{enumerate}

\section{Rollback}
Rollback is intentionally simple:
\begin{enumerate}[leftmargin=1.5em]
    \item list prior release directories,
    \item repoint \code{current} to the previous known-good release,
    \item restart the systemd service,
    \item re-run the local and public health checks.
\end{enumerate}

\section{Operational Smoke Checks}
After every deploy, run:
\begin{lstlisting}
curl -fsS http://127.0.0.1:8008/health
curl -fsS https://your-hostname.example.com/health
\end{lstlisting}
Then confirm interactive behavior by logging in and opening:
\begin{itemize}[leftmargin=1.5em]
    \item \code{/dashboard},
    \item \path{/trade/equity_cliquet}, and
    \item \code{/quantlib-data-model}.
\end{itemize}

\section{TLS}
For HTTPS, use Apache with a real certificate. On a public hostname, \code{certbot} is the lowest-friction path:
\begin{lstlisting}
sudo apt-get install -y certbot python3-certbot-apache
sudo certbot --apache -d your-hostname.example.com
\end{lstlisting}

\section{Notes}
This deployment model intentionally stays small and inexpensive:
\begin{itemize}[leftmargin=1.5em]
    \item one droplet,
    \item local filesystem storage,
    \item no database,
    \item no extra worker host, and
    \item no paid external platform dependencies beyond the droplet itself.
\end{itemize}

\end{document}
